\documentclass[letterpaper,11pt]{article}
\def\theme{MidnightBlue}
\usepackage[style=ieee,url=false,doi=false,eprint=true,maxbibnames=99,sorting=none,dashed=false]{biblatex}
\usepackage{cv}

\newcommand{\prl}{Phys.\ Rev.\ Lett.} 
\newcommand{\prd}{Phys.\ Rev.\ D}
\newcommand{\mnras}{Mon.\ Not.\ R.\ Astron.\ Soc.}
\newcommand{\apj}{Astrophys.\ J.}%
\newcommand{\apjl}{Astrophys.\ J.\ Lett.}
\newcommand{\apjs}{Astrophys.\ J.\ Suppl.}
\newcommand{\prx}{Phys.\ Rev.\ X}

\boldname{Vajpeyi}{Avi}{A.}
\addbibresource{papers/papers.bib}
\hypersetup{pdftitle={Avi Vajpeyi - CV - Anthropic STEM Fellows},pdfauthor={Avi Vajpeyi}}

\begin{document}
\headinginline{Avi Vajpeyi}{
  \email{avi.vajpeyi@auckland.ac.nz}\\
  \website{avivajpeyi.github.io}\\
  \website{github.com/avivajpeyi}
}

\section{Research Profile}
Astrostatistician developing Bayesian methods and research software for noisy,
correlated time series. I work from mathematical models through implementation
and simulation studies, checking when assumptions hold and where approximations
fail. PhD in astrophysics; background in computer science and physics.

\smallskip
\textbf{US permanent resident; no visa sponsorship required.}

\import{sections/}{positions.tex}

\section{Selected Research and Scientific Software}
\denseouterlist{
\entrymid[\textbullet]
{\textbf{Bayesian spectral estimation} --- \link{https://github.com/nz-gravity/LogPSplinePSD}{LogPSplinePSD}}{}
{Develop P-spline models and JAX/NumPyro inference for correlated and
time-varying detector noise. Implement variational initialisation and HMC/NUTS;
test recovery of known spectral matrices and uncertainty calibration on
simulations. First-author multivariate study tests the diagonal-noise
approximation against asymmetric LISA noise; co-authored VI study compares
posterior approximations with HMC.}

\entrymid[\textbullet]
{\textbf{Differentiable time--frequency analysis} --- \link{https://github.com/pywavelet/wdm_transform}{wdm-transform}}{}
{First author and software developer of an explicit Wilson--Daubechies--Meyer
transform for gravitational-wave analysis. Connect mathematical derivations
with NumPy/JAX implementations and numerical checks of reconstruction,
orthogonality and coefficient covariance, supporting inference with
time-localised noise models.}

\entrymid[\textbullet]
{\textbf{Surrogates for expensive simulations} --- \link{https://github.com/COMPAS-Surrogate/lnl_surrogate}{COMPAS LnL Surrogate}}{}
{Built a likelihood-surrogate package linking COMPAS population simulations
to Gaussian-process models and Bayesian optimisation. Implemented active
learning with exploratory and exploitative acquisition functions, with mock
likelihood examples for checking surrogate predictions and uncertainty.}

\entrymid[\textbullet]
{\textbf{Generative waveform models} --- \link{https://github.com/starccato/starccato_jax}{Starccato-JAX}}{}
{Co-develop JAX waveform generators for core-collapse supernova signals and
detector glitches. The package provides pretrained decoders and reproducible
latent-space sampling for waveform studies; guide related student research
on supernova detection with normalising flows.}
}

\section{Tools and Working Practice}
Python, NumPy/SciPy, JAX, NumPyro, ArviZ; TensorFlow, Gaussian processes,
variational inference, HMC, nested sampling; Git, testing and computing-cluster
workflows. Use Claude Code and other coding agents daily to implement
mathematical models and run cluster jobs, reviewing code and validating results.

\newpage
\section{Research Software Contributions}
\denseouterlist{
\entrymid[\textbullet]
{\textbf{\link{https://arxiv.org/abs/1909.11873}{parallel\_bilby}, bilby / bilby\_pipe and asimov}}{}
{Co-developed parallel\_bilby for computationally demanding LVK parameter
estimation; led its ADACS-funded upgrade for the fourth observing run.
Contributed maintenance to bilby/bilby\_pipe and to asimov's coordination of
parameter-estimation workflows on computing clusters.}
}

\section{Industry Experience}
\denseouterlist{
\entrymid[\textbullet]{\textbf{Software Consultant, Quaisr}}{2022--2023}
{Built data-agnostic ML and active-learning blocks linking TensorFlow backends
to visual workflows.}
\entrymid[\textbullet]{\textbf{Software Engineer, Bamboo Learning}}{2018--2019}
{Developed NLP components and maths/listening-comprehension content for Alexa.}
\entrymid[\textbullet]{\textbf{Machine Learning Consultant, MOAD}}{2018}
{Deployed CNN image classifiers and delivered hands-on machine-learning tutorials.}
}

\section{Education}
\textbf{PhD, Astrophysics}, Monash University, 2019--2022.
Thesis: \link{https://doi.org/10.26180/21800807}{Astrophysical Inference in the Data Driven Era}.\\
\textbf{BA, Computer Science and Physics}, College of Wooster, 2014--2018.
GPA 4.0/4.0; valedictorian. Undergraduate thesis on GPU granular-flow simulations.

\section{Selected Publications}
{\small
\nocite{Liu:2026:2608.22245,Vajpeyi:2026:arXiva,Vajpeyi:2026:arXivb,Zahraoui:2026:PhRvD,Liu:2025:PhRvD,Smith:2020:MNRAS}
\printbibliography[heading=none]
}

\section{Communication and Mentoring}
Provide day-to-day postgraduate supervision in Bayesian noise estimation and waveform
modelling. Independently mentored software-engineering student teams at Monash.
Delivered a LISA research talk and a separate Bayesian parameter-estimation
training lecture in Beijing (August--September 2026); BayesComp invited session
on Bayesian computation in astrophysics (2025).
\end{document}
